\xiaosi\song

本课程结课大作业的完成，离不开各位老师和同学的帮助与支持，在此表示衷心的感谢。

首先，感谢卢炳先老师一学期来的悉心教导。卢老师在课堂上深入浅出地讲解了机器学习的核心理论与前沿方法，从监督学习、无监督学习到深度学习、强化学习，构建了完整的知识体系。这些系统化的知识为本大作业的选题和实施奠定了坚实基础。特别感谢卢老师在课程中对神经网络架构、可解释性等主题的讲解，直接启发了本文"基于神经电路策略的可解释自动驾驶规划"这一研究方向的形成。

其次，感谢课程中一起学习讨论的同学们。在课后讨论与实验交流中，同学们提出的问题和建议让我从多个角度重新审视本文的方法设计，改进了许多细节。

感谢神经电路策略（NCP）和液态时间常数网络（LTC）原作者Lechner、Hasani等人开源的高质量代码库，以及Highway-env仿真平台的贡献者们，为本大作业的实验工作提供了便利。

由于时间和能力有限，本大作业仍存在一些不足之处——例如消融实验的种子数量较少、搜索环境的覆盖不够全面等。恳请卢老师批评指正。

\begin{flushright}
\song\xiaosi
孙家恒\\
二〇二六年五月
\end{flushright}
